% THE SECOND $ OF A DISPLAY'S $$ IS FETCHED WITH EXPANSION.
%
% When a display's first closing $ has been read, the reference looks for
% the second one the way it looks for anything else -- expanding as it
% goes. So a macro standing for it will do, and \expandafter may stand
% between the two:
%
%   \def\dollar{$}   ... $$y$\dollar        closes cleanly, no complaint
%                    ... $$y$\relax         "Display math should end with $$"
%                    ... $$y$\expandafter$\csname dollar\endcsname
%
% HSTeX read the token raw, so the first of those drew a complaint.
%
% The mode there is the list the display will join, not the display's --
% the formula's own list is gone by then. Inside a \vbox that is internal
% vertical mode, which is what \tracingcommands names for the
% \expandafter of the third case.
%
% What \expandafter expands is traced like anything the expander expands,
% so {\csname} is drawn there too.
%
% An equation NUMBER's closing $ is different: the list that ends there
% was part of the display, and the display is still standing, so the mode
% for the second $ is display math mode.
\catcode`\{=1 \catcode`\}=2 \catcode`\$=3 \catcode`\#=6
\tracingonline=1 \hsize=200pt \parindent=0pt \hbadness=10000 \hfuzz=100pt
\def\dollar{$}
\message{[A the second $ comes from a macro]}
\tracingcommands=2
\setbox0=\vbox{\noindent x$$y$\dollar z\par}
\tracingcommands=0
\message{[B the second $ does not come at all]}
\tracingcommands=2
\setbox0=\vbox{\noindent x$$y$\relax z\par}
\tracingcommands=0
\message{[C the second $ comes from \expandafter]}
\tracingcommands=2
\setbox0=\vbox{\noindent x$$y$\expandafter$\csname dollar\endcsname z\par}
\tracingcommands=0
\message{[D an equation number closed with expansion between the shifts]}
\tracingcommands=2
\setbox0=\vbox{\noindent x$$y\eqno z$\expandafter$\csname dollar\endcsname w\par}
\tracingcommands=0
\message{[E an equation number closed plainly]}
\tracingcommands=2
\setbox0=\vbox{\noindent x$$y\eqno z$$w\par}
\tracingcommands=0
\message{[done]}
\end
